% ============================================================================
% Lab Chapter for Part V — Practicals 1–8
% ============================================================================
\chapter{Part V --- Hands-On Practicals}
\label{ch:lab5}

\bytetip[fun]{Spreadsheet time! These practicals will take you from entering your first number to creating charts that would impress your maths teacher.}

\begin{labactivity}{1}{Exploring the Excel Interface}{10 min}
\youwillneed{\faDesktop\ Computer \quad \faClock\ 10 min \quad \faFileExcel\ MS Excel}
\textbf{Objective:} Identify all parts of the Excel window.\\[4pt]
\textbf{Steps:}
\begin{enumerate}[label=\protect\stepcircle{\arabic*}, leftmargin=2cm, itemsep=4pt]
  \item Open Excel and create a new blank workbook.
  \item Click on cell A1. Notice the cell address in the Name Box.
  \item Click on cell D5. What does the Name Box show now?
  \item Look at the Formula Bar --- what does it display?
  \item Find the Sheet Tabs at the bottom. How many sheets are there?
\end{enumerate}
\textbf{Expected Outcome:} Students identify Name Box, Formula Bar, and Sheet Tabs.
\welldone
\end{labactivity}

\begin{labactivity}{2}{Entering and Formatting Data}{15 min}
\youwillneed{\faDesktop\ Computer \quad \faClock\ 15 min \quad \faFileExcel\ MS Excel}
\textbf{Objective:} Enter data and apply basic formatting.\\[4pt]
\textbf{Steps:}
\begin{enumerate}[label=\protect\stepcircle{\arabic*}, leftmargin=2cm, itemsep=4pt]
  \item In A1, type ``Student Name.'' In B1, type ``Maths.'' In C1, type ``Science.''
  \item Enter 5 students' names and marks in the rows below.
  \item Bold the header row. Change the header background to blue with white text.
  \item Apply borders to all data cells (Home $\rightarrow$ Borders).
  \item Adjust column widths by double-clicking the column border.
  \item Save as ``StudentMarks.xlsx.''
\end{enumerate}
\textbf{Expected Outcome:} A formatted table with 5 students' marks.
\welldone
\end{labactivity}

\begin{labactivity}{3}{Using AutoFill}{10 min}
\youwillneed{\faDesktop\ Computer \quad \faClock\ 10 min \quad \faFileExcel\ MS Excel}
\textbf{Objective:} Use AutoFill to create data series.\\[4pt]
\textbf{Steps:}
\begin{enumerate}[label=\protect\stepcircle{\arabic*}, leftmargin=2cm, itemsep=4pt]
  \item In A1 type ``Monday.'' Use AutoFill to fill A1:A7 with days of the week.
  \item In B1 type ``January.'' Use AutoFill to fill B1:B12 with months.
  \item In C1 type 2, in C2 type 4. Select both and AutoFill to C10 (even numbers).
  \item In D1 type ``Quarter 1.'' AutoFill down to D4.
  \item Observe how Excel recognises each pattern!
\end{enumerate}
\textbf{Expected Outcome:} Students fill 4 different series using AutoFill.
\welldone
\end{labactivity}

\begin{labactivity}{4}{Writing Basic Formulas}{15 min}
\youwillneed{\faDesktop\ Computer \quad \faClock\ 15 min \quad \faFileExcel\ MS Excel}
\textbf{Objective:} Create formulas using +, -, *, /.\\[4pt]
\textbf{Steps:}
\begin{enumerate}[label=\protect\stepcircle{\arabic*}, leftmargin=2cm, itemsep=4pt]
  \item Open ``StudentMarks.xlsx'' from Practical 2.
  \item In D1 type ``Total.'' In D2 type \texttt{=B2+C2} and press Enter.
  \item Use the fill handle to copy the formula down for all 5 students.
  \item In E1 type ``Average.'' In E2 type \texttt{=(B2+C2)/2} and fill down.
  \item Click on D2 and check the formula in the Formula Bar.
  \item Save the file.
\end{enumerate}
\textbf{Expected Outcome:} Students calculate totals and averages using formulas.
\welldone
\end{labactivity}

\begin{labactivity}{5}{Using Functions (SUM, AVERAGE, MAX, MIN)}{15 min}
\youwillneed{\faDesktop\ Computer \quad \faClock\ 15 min \quad \faFileExcel\ MS Excel}
\textbf{Objective:} Use built-in functions to analyse data.\\[4pt]
\textbf{Steps:}
\begin{enumerate}[label=\protect\stepcircle{\arabic*}, leftmargin=2cm, itemsep=4pt]
  \item Continue with ``StudentMarks.xlsx.''
  \item In B8 type ``Sum:'' In C8 type \texttt{=SUM(B2:B6)}. What result do you get?
  \item In B9 type ``Average:'' In C9 type \texttt{=AVERAGE(B2:B6)}.
  \item In B10 type ``Highest:'' In C10 type \texttt{=MAX(B2:B6)}.
  \item In B11 type ``Lowest:'' In C11 type \texttt{=MIN(B2:B6)}.
  \item Repeat for the Science column (column C).
  \item Save.
\end{enumerate}
\textbf{Expected Outcome:} Students use 4 functions correctly.
\welldone
\end{labactivity}

\begin{labactivity}{6}{Number Formatting}{10 min}
\youwillneed{\faDesktop\ Computer \quad \faClock\ 10 min \quad \faFileExcel\ MS Excel}
\textbf{Objective:} Apply different number formats to cells.\\[4pt]
\textbf{Steps:}
\begin{enumerate}[label=\protect\stepcircle{\arabic*}, leftmargin=2cm, itemsep=4pt]
  \item Create a new sheet. Enter prices of 5 items in column B.
  \item Select B1:B5. Apply \textbf{Currency} format (\key{Ctrl}+\key{Shift}+\key{\$}).
  \item In C1:C5 enter percentages (e.g., 0.85). Apply \textbf{Percentage} format.
  \item In D1:D5 enter dates. Apply a date format (Format Cells $\rightarrow$ Date).
  \item Compare how the same number looks in different formats.
\end{enumerate}
\textbf{Expected Outcome:} Students apply 3 different number formats.
\welldone
\end{labactivity}

\begin{labactivity}{7}{Creating Charts}{15 min}
\youwillneed{\faDesktop\ Computer \quad \faClock\ 15 min \quad \faFileExcel\ MS Excel}
\textbf{Objective:} Create bar and pie charts from data.\\[4pt]
\textbf{Steps:}
\begin{enumerate}[label=\protect\stepcircle{\arabic*}, leftmargin=2cm, itemsep=4pt]
  \item Open ``StudentMarks.xlsx.''
  \item Select cells A1:C6 (names and marks for both subjects).
  \item Insert $\rightarrow$ Charts $\rightarrow$ \textbf{Clustered Bar Chart}.
  \item Add a chart title: ``Student Marks Comparison.''
  \item Now select one student's marks to create a \textbf{Pie Chart}.
  \item Add data labels to the pie chart (Chart Design $\rightarrow$ Add Chart Element).
  \item Save.
\end{enumerate}
\textbf{Expected Outcome:} A bar chart and a pie chart from the same data.
\welldone
\end{labactivity}

\begin{labactivity}{8}{Sorting and Complete Spreadsheet}{15 min}
\youwillneed{\faDesktop\ Computer \quad \faClock\ 15 min \quad \faFileExcel\ MS Excel}
\textbf{Objective:} Sort data and create a polished spreadsheet.\\[4pt]
\textbf{Steps:}
\begin{enumerate}[label=\protect\stepcircle{\arabic*}, leftmargin=2cm, itemsep=4pt]
  \item Open ``StudentMarks.xlsx.''
  \item Click anywhere in the Total column. Data tab $\rightarrow$ Sort Largest to Smallest.
  \item Observe how all rows rearrange together.
  \item Add conditional formatting: select marks $\rightarrow$ Home $\rightarrow$ Conditional Formatting $\rightarrow$ Colour Scales.
  \item Add a merged title row at the top: ``Class 8 --- Term 1 Results.''
  \item Save and check Print Preview.
\end{enumerate}
\textbf{Expected Outcome:} A professional sorted and formatted spreadsheet.
\welldone
\end{labactivity}

\vspace{12pt}
\begin{center}
  \qrcode[height=2cm]{https://example.com/module5}\\[6pt]
  {\small\sffamily\textcolor{NavyBlue}{Scan for extra Part V resources and activities!}}
\end{center}
